\section{组织治理：如何防止 benchmark 被 ``刷分工程'' 劫持}

\subsection{评测治理不只是技术问题}
Sida 在问答里提到，噪声问题最终会变成组织问题：如果团队激励只看 leaderboard，系统自然会朝 ``短期可刷分'' 方向演化。治理目标是让激励函数和真实能力一致，而不是让发布节奏被单一分数绑架。

\begin{importantbox}{评测治理的三条红线}
\begin{enumerate}
    \item 禁止只报 best run，不报完整 run distribution；
    \item 禁止跨协议比较但仍给出确定性排序；
    \item 禁止把私有调参结果当作通用能力结论对外发布。
\end{enumerate}
\end{importantbox}

\subsection{抗游戏化（anti-gaming）机制}
为了降低 benchmark gaming，可以采取 challenge split、协议轮换、隐藏测试集、以及周期性人工审计。课程强调，治理不是为了变慢，而是为了让 ``快'' 不以牺牲真实性为代价。

\begin{table}[H]
\centering
\begin{tabular}{p{3.1cm}p{3.8cm}p{3.8cm}}
\toprule
\textbf{机制} & \textbf{解决的问题} & \textbf{代价} \\
\midrule
Hidden Test Split & 防止公开集过拟合 & 复现成本上升 \\
Protocol Rotation & 防止模板投机 & 历史可比性下降 \\
Human Audit & 校准 judge 偏差 & 人力成本较高 \\
Release Gating & 把统计门槛写入流程 & 发布速度变慢 \\
\bottomrule
\end{tabular}
\caption{治理机制是在 ``真实性'' 与 ``效率'' 间做工程折中}
\end{table}

\subsection{面向 2026 的评测系统形态}
如果把本讲结论外推到下一阶段，团队需要的不只是 benchmark harness，而是 benchmark platform：可版本化协议、可追踪日志、可配置统计报告、可插拔 judge 和可审计发布流程。只有这样，benchmark 才能成为研发基础设施，而不是一次性演示工具。

\begin{warningbox}{没有治理，噪声会反向塑造研究方向}
当团队把 ``噪声中奖'' 误当成 ``能力突破''，研究资源会被系统性错配，长期会拖慢真实进展。
\end{warningbox}

\subsection{本章小结}
Benchmark 可靠性最终由组织机制保证。技术框架给出测量能力，治理机制保证测量结果不会被短期激励扭曲。

